\chapter{Électricité, matières premières et connectivité des bourses de paris}\label{ch:nw:power-commodities-and-betting-connectivity}

Le marché infrajournalier européen de l'électricité transfrontalier apparie les ordres de vingt-cinq pays dans un seul \omterm{def:nw:power-commodities-and-betting-connectivity:sob}{carnet d'ordres partagé}, et
répartit la capacité des interconnecteurs à chaque transaction, premier arrivé, premier servi. Lorsque la capacité à travers une frontière est limitée,
le dernier mégawatt revient à celui qui atteint le carnet en premier. Dans le modèle du chapitre, un participant situé à \qty{10}{\milli\second} du
carnet, face à un rival à cinq, obtient le dernier mégawatt dans 4.5~\% des contestations~; à \qty{2}{\milli\second}, dans 98.5~\%.

Le livre~3 décrivait l'électricité infrajournalière, le couplage des marchés, les zones de dépôt des offres, l'équilibrage et les nominations de programme, ainsi que les bourses de paris.
Ce chapitre traite de la manière dont les firmes s'y connectent~: le \omterm{def:nw:power-commodities-and-betting-connectivity:sob}{carnet d'ordres partagé} et les systèmes de trading locaux placés devant lui, les
systèmes de programmation des gestionnaires de réseau, les bourses de matières premières de la carte du chapitre~11, et les interfaces et
les frais d'une bourse de paris. La course à la capacité est une simulation étiquetée~; les règles sont celles des opérateurs eux-mêmes.

\section{Interfaces de l'électricité infrajournalière et carnet d'ordres partagé}

\begin{definition}[Carnet d'ordres partagé, allocation continue implicite]\label{def:nw:power-commodities-and-betting-connectivity:sob}
Un \emph{carnet d'ordres partagé}\index{carnet d'ordres partagé} est un seul carnet d'ordres que plusieurs opérateurs de marché alimentent avec les ordres de leurs
participants, de sorte qu'un ordre saisi auprès d'un opérateur puisse être apparié à un ordre saisi auprès d'un autre.
L'\emph{allocation continue implicite}\index{allocation continue implicite} est l'allocation de la capacité de transport transfrontalière
à l'instant où deux ordres situés dans des zones de dépôt des offres différentes s'apparient, la capacité et l'énergie étant tarifées ensemble, et la
capacité utilisée par la transaction est aussitôt retirée de ce que la suivante peut utiliser.
\end{definition}

Le couplage infrajournalier unique (SIDC) fait fonctionner un seul \omterm{def:nw:power-commodities-and-betting-connectivity:sob}{carnet d'ordres partagé}, un
module de gestion des capacités, qui recense les capacités interzonales libérées par les gestionnaires de réseau, et un module d'acheminement
(l'\cref{dat:nw:power-commodities-and-betting-connectivity:sidc}). Un participant ne se connecte pas directement au carnet partagé~: il
passe par le système de trading local de son opérateur de marché, qui transmet les ordres au carnet et lui montre les ordres des
participants des autres opérateurs, dans la mesure où la capacité le permet. Au sein d'un niveau de prix, les ordres sont classés par leur heure d'enregistrement~; à travers une
frontière, un appariement n'est possible que tant qu'il reste de la capacité, et chaque appariement la consomme.

\begin{dated}{2026-09}{Le couplage infrajournalier unique, selon l'ENTSO-E et un opérateur de marché}\label{dat:nw:power-commodities-and-betting-connectivity:sidc}
\omterm{def:nw:power-commodities-and-betting-connectivity:sob}{Carnet d'ordres partagé}, module de gestion des capacités et module d'acheminement~; «~la transaction est conclue selon le principe du premier arrivé, premier servi
où le prix d'achat le plus élevé et le prix de vente le plus bas sont servis en premier~»~; allocation implicite (explicite seulement aux frontières franco-allemandes
et croato-slovènes)~; 25 pays couplés depuis la quatrième vague (novembre 2022)~; l'heure de fermeture du guichet interzonal doit passer
de 60 à 30 minutes avant la livraison. Les règles de Nord Pool (juin 2024)~: prix, puis heure d'enregistrement~; allocation interzonale sur le marché continu
suspendue pendant les enchères infrajournalières sur les mêmes frontières.
\end{dated}

Une frontière à capacité limitée est une course. On imagine un écart de prix qui s'ouvre entre deux zones, avec encore de la place sur l'interconnecteur pour un
ordre de plus~: chaque participant qui le voit envoie un ordre, et le premier à atteindre le carnet partagé obtient la capacité. Le modèle du chapitre
compte cinq rivaux dont les latences vers le carnet suivent une loi lognormale centrée autour de 5, 10, 20, 40 et \qty{80}{\milli\second}, avec une dispersion de 0.3, et
fait varier sa propre médiane de latence (toutes les hypothèses~: l'hébergement du carnet et la vitesse de chaque système local ne sont pas publiés).

\omcode{code/firm/powerlink/firm_powerlink.py}{35}{41}{La course à la capacité~: les premiers arrivés obtiennent la capacité rare.}\label{lst:nw:power-commodities-and-betting-connectivity:race}

\begin{proposition}[La rareté rend la latence tout ou rien]\label{prop:nw:power-commodities-and-betting-connectivity:scarce}
Lorsque la capacité restante équivaut à $s$ ordres et que $n$ participants se disputent cette capacité, un participant en obtient si et seulement s'il fait partie des
$s$ premiers arrivants~; avec une seule place, sa part est la probabilité d'être le plus rapide, qui passe de un à zéro sur une plage de
latences fixée par la dispersion du rival le plus rapide~; avec davantage de places, la chute se déplace vers le $s$-ième rival le plus rapide.
\end{proposition}

\begin{proof}
L'allocation se fait selon l'ordre d'arrivée, donc un participant gagne exactement lorsque moins de $s$ autres arrivent avant lui. Avec une seule place, cela revient à
l'événement où sa latence est inférieure à celle de chaque rival~; cette probabilité décroît avec sa propre latence et est dominée par celle du rival le plus rapide.
\end{proof}

\begin{omfigure}
\begin{tikzpicture}
\begin{axis}[width=11cm,height=5.8cm,xmode=log,xmin=1,xmax=100,ymin=0,ymax=1.05,
  xlabel={notre latence médiane vers le carnet d'ordres partagé (ms, log)},ylabel={part des contestations remportées},
  tick label style={font=\footnotesize},label style={font=\small},grid=major,grid style={gray!25},
  legend style={font=\footnotesize,at={(0.5,-0.3)},anchor=north,legend columns=2},table/col sep=comma]
\addplot[omThm,thick,mark=*,mark size=1.4pt] table[x=latency_ms,y=one]{figdata/networks/26-power-commodities-and-betting-connectivity/race.csv};
\addplot[omDef,thick,mark=square*,mark size=1.4pt] table[x=latency_ms,y=three]{figdata/networks/26-power-commodities-and-betting-connectivity/race.csv};
\legend{capacité restante pour un ordre,capacité restante pour trois ordres}
\end{axis}
\end{tikzpicture}

\omcaption{Simulation~: la part de la capacité transfrontalière rare remportée par la firme en fonction de sa latence médiane vers le \omterm{def:nw:power-commodities-and-betting-connectivity:sob}{carnet d'ordres partagé},
face à cinq rivaux à 5, 10, 20, 40 et \qty{80}{\milli\second}. Avec une capacité restante pour un ordre, la part passe de 98.5~\% à
\qty{2}{\milli\second} à 4.5~\% à \qty{10}{\milli\second}~; avec trois ordres, elle se maintient à 95.4~\% à \qty{10}{\milli\second}. Données~:
\texttt{fig\_power.py}, \texttt{nw\_power.share\_curve()}.}\label{fig:nw:power-commodities-and-betting-connectivity:race}
\end{omfigure}

La valeur en découle. Si l'écart de prix que capte la capacité atteint en moyenne 12~euros par mégawattheure (une hypothèse), une firme située à
\qty{2}{\milli\second} espère 11.82~euros par mégawattheure offert dans de telles contestations, et une firme à \qty{10}{\milli\second} 0.54~euro~:
toute la valeur d'une frontière à capacité limitée revient au participant le plus rapide, et le second n'en touche presque rien. Avec une capacité restante pour trois
ordres, la même firme à \qty{10}{\milli\second} obtient 95.4~\% des contestations. La rareté, et non le marché, transforme un desk de trading d'électricité infrajournalière
en une activité pilotée par la latence, et seulement sur les frontières et aux heures où la capacité devient contraignante.

L'heure de fermeture du guichet borne la course dans le temps~: le trading interzonal pour une période de livraison s'arrête 60 minutes avant celle-ci (30 minutes après le
changement prévu), et pendant les enchères infrajournalières, l'allocation interzonale du marché continu s'arrête sur les frontières mises aux enchères. L'horloge
d'un desk de trading comporte donc trois types de moments~: le trading continu avec capacité, les fenêtres d'enchères sans capacité, et la fermeture du guichet.

\section{Systèmes de programmation des gestionnaires de réseau}

Après le trading, les positions deviennent des programmes~: un responsable d'équilibre soumet à son gestionnaire de réseau une nomination de programme pour l'énergie qu'il va
injecter, soutirer et échanger, période par période, avant les échéances du gestionnaire. Ces échanges ne sont pas des messages de trading mais
des documents, transmis par l'Energy Communication Platform de l'ENTSO-E, qui assure une livraison fiable et sécurisée par AMQP, le service web MADES
(IEC 62325-504) ou par des dossiers partagés. Ici, la latence se compte en minutes, non en millisecondes~; l'essentiel est que chaque transaction ait un programme, que
les documents soient acquittés, et qu'une livraison échouée soit repérée avant l'échéance, soit la même discipline de rapprochement que celle du
drop copy (chapitre~23), à une cadence plus lente.

\section{Connectivité des bourses de matières premières}

Les bourses de contrats à terme sur matières premières sont des bourses comme les autres, sur la carte du chapitre~11~: l'appariement d'ICE Futures Europe à Basildon, les
contrats européens sur l'énergie et les métaux, sur les mêmes campus continentaux que les contrats financiers. Leur connectivité est celle du chapitre~9~:
les \omterm{def:nw:colocation-products-and-how-they-are-sold:mmr}{interconnexions} (cross-connects), la \omterm{def:nw:colocation-products-and-how-they-are-sold:colo}{colocation}, les feeds et les passerelles, avec les questions d'équité du chapitre~23. Ce qui change pour un desk d'électricité ou de gaz est
la combinaison~: une position sur un contrat à terme en bourse, une position infrajournalière passant par un système de trading local et par le carnet partagé, et
des programmes adressés à un gestionnaire de réseau, trois connexions sur trois horloges pour une seule exposition physique.

\section{Interfaces des bourses de paris et leurs frais de données}

Une bourse de paris vend l'accès à ses carnets d'ordres par une interface de programmation applicative, et en fixe le prix selon ce que le client
lui demande. La documentation de Betfair est explicite (l'\cref{dat:nw:power-commodities-and-betting-connectivity:betfair})~: les requêtes de données sont
limitées par un poids, 200 points par requête, chaque marché d'une requête coûtant le poids de la projection demandée~; les transactions
au-delà d'un seuil horaire sont facturées~; et la bourse recommande son interface de streaming plutôt que le polling.

\begin{definition}[Frais de transaction]\label{def:nw:power-commodities-and-betting-connectivity:txn}
Les \emph{frais de transaction}\index{frais de transaction} sont une commission de la plateforme sur le nombre d'actions sur ordres qu'un participant envoie au-delà d'un
seuil gratuit par période, en comptant les placements et les actions échouées, quelle qu'en soit l'issue, pour tarifer la charge
qu'imposent les re-cotations fréquentes à la plateforme.
\end{definition}

\begin{dated}{2026-09}{Limites de requêtes et frais d'une bourse de paris, d'après sa documentation développeur}\label{dat:nw:power-commodities-and-betting-connectivity:betfair}
API Betfair Exchange~: le poids d'une requête (le poids de la projection multiplié par le nombre de marchés) ne doit pas dépasser 200 points~; les meilleurs cours vendeurs
pèsent 5 par marché, tous les cours vendeurs avec volumes échangés 32. Des \omterm{def:nw:power-commodities-and-betting-connectivity:txn}{frais de transaction} peuvent s'appliquer au-delà de 5\,000 transactions placées ou échouées par heure~;
un placement et son annulation réussie ne comptent qu'une fois. Clé d'application en production~: des frais uniques de 499~GBP. «~Utiliser l'API Stream plutôt
que le polling chaque fois que possible.~»
\end{dated}

\omcode{code/firm/powerlink/firm_powerlink.py}{53}{67}{Arithmétique du polling à partir des poids publiés, et transactions au-delà du seuil horaire.}\label{lst:nw:power-commodities-and-betting-connectivity:poll}

Faire du polling sur 1\,000 marchés cinq fois par seconde pour les meilleurs cours vendeurs suppose 40 marchés par requête, 25 requêtes par cycle et 125 requêtes par seconde~; pour
la totalité des niveaux de prix avec volumes échangés, 6 marchés par requête et 835 requêtes par seconde. Dans les deux cas, les données accusent en moyenne
un retard de \qty{110}{\milli\second} avec un aller-retour de \qty{10}{\milli\second}, soit la moitié de l'intervalle de polling plus le trajet. Un flux envoie chaque
changement une seule fois, au moment où il survient. Et une stratégie qui re-cote 30 marchés quatre fois par minute envoie 7\,200 transactions par heure, 2\,200 au-delà
du seuil~: sur une plateforme qui facture à l'action, les habitudes d'annulation-remplacement d'un teneur de marché en bourse constituent un poste de coût.

\begin{method}[Connecter un desk d'électricité ou de paris]\label{met:nw:power-commodities-and-betting-connectivity:plan}
\begin{enumerate}
\item Associer chaque marché à son horloge~: en continu avec capacité, enchères, heures de fermeture du guichet, échéances de nomination de programme.
\item Pour l'électricité infrajournalière, mesurer la latence jusqu'au \omterm{def:nw:power-commodities-and-betting-connectivity:sob}{carnet d'ordres partagé} à travers chaque système de négociation local accessible, et repérer les frontières
  et les heures où la capacité est contraignante~; c'est là que la latence rapporte.
\item Transmettre les programmes via la plateforme des gestionnaires de réseau avec accusés de réception, et rapprocher les transactions des programmes avant chaque
  échéance.
\item Pour les bourses de paris, utiliser le streaming plutôt que le polling, et compter les transactions par rapport au seuil facturé.
\item Mettre les contrats à terme sur matières premières au niveau du chapitre~9~: \omterm{def:nw:colocation-products-and-how-they-are-sold:colo}{colocation}, flux de données, passerelles.
\end{enumerate}
\end{method}

\begin{omfigure}
\begin{tikzpicture}[font=\small,
  box/.style={draw=omDef,fill=omDef!10,rounded corners,align=center,minimum height=9mm,inner sep=3pt},
  ops/.style={draw=omThm,fill=omThm!10,rounded corners,align=center,minimum height=9mm,inner sep=3pt}]
\node[box] (a) at (0,1.6) {participant\\zone A};
\node[box] (b) at (0,-0.2) {participant\\zone B};
\node[ops] (la) at (3.6,1.6) {système de négociation\\local (opérateur 1)};
\node[ops] (lb) at (3.6,-0.2) {système de négociation\\local (opérateur 2)};
\node[ops,minimum width=3.2cm] (sob) at (8.2,0.7) {carnet d'ordres partagé\\gestion de la capacité};
\node[box] (tso) at (8.2,-1.6) {gestionnaires de réseau\\(capacités)};
\draw[->,omDef] (a) -- (la); \draw[->,omDef] (b) -- (lb);
\draw[->,omThm,thick] (la) -- (sob); \draw[->,omThm,thick] (lb) -- (sob);
\draw[->,omDef,dashed] (tso) -- node[right,font=\footnotesize] {capacité inter-zones} (sob);
\end{tikzpicture}

\omcaption{Le couplage infrajournalier unique~: les participants négocient via le système de négociation local de leur opérateur~; les ordres se rencontrent dans un seul \omterm{def:nw:power-commodities-and-betting-connectivity:sob}{carnet d'ordres partagé},
et chaque appariement transfrontalier consomme la capacité libérée par les gestionnaires de réseau. Schéma de principe.}\label{fig:nw:power-commodities-and-betting-connectivity:sidc}
\end{omfigure}

\section{Tutoriel~: le dernier mégawatt}\label{tut:nw:power-commodities-and-betting-connectivity:tut}

\begin{tutorial}
\textbf{Objectif.} Simuler une course à la capacité transfrontalière rare à plusieurs latences, et chiffrer le coût du polling d'une bourse de paris par rapport au streaming.
\textbf{État final~:} la \cref{fig:nw:power-commodities-and-betting-connectivity:race} et les chiffres de polling de la section~4.
\begin{tutsteps}
\item \textbf{Course.} \texttt{firm\_powerlink.capacity\_race} (le \cref{lst:nw:power-commodities-and-betting-connectivity:race}) avec la
  firme et cinq concurrents~; \texttt{nw\_power.share\_curve} pour un ordre et pour trois ordres de capacité.
\item \textbf{Valorisation.} \texttt{nw\_power.value} au spread moyen retenu.
\item \textbf{Polling.} \texttt{load\_weights}, \texttt{poll\_requests\_per\_s} et \texttt{charged\_transactions}
  (le \cref{lst:nw:power-commodities-and-betting-connectivity:poll}).
\end{tutsteps}
\textbf{Pour aller plus loin.} Faire dépendre la taille de l'écart de prix de la vitesse à laquelle il se referme, et ajouter le temps de traitement propre
au module de gestion de la capacité, comme délai commun.
\end{tutorial}

\section{Construction~: le modèle de connectivité électricité et paris}\label{bld:nw:power-commodities-and-betting-connectivity:powerlink}

\begin{build}
\bfield{Rôle} La course à la capacité, l'horloge de fermeture du guichet et les budgets de requêtes et de transactions d'une bourse de paris~: les lignes électricité,
matières premières et paris du plan du chapitre~29.
\bfield{Interface} \texttt{firm\_powerlink}~: \texttt{Participant}, \texttt{capacity\_race}, \texttt{gate\_closure},
\texttt{load\_weights}, \texttt{markets\_per\_request}, \texttt{poll\_requests\_per\_s}, \texttt{poll\_staleness\_ms},
\texttt{charged\_transactions}.
\bfield{Règles} Les poids et les seuils sont des lignes datées avec sources~; les latences sont des hypothèses énoncées~; la course est une simulation étiquetée.
\bfield{Tests d'acceptation} \texttt{code/firm/powerlink/tests/}~: des participants égaux se partagent la capacité à parts égales, des latences fixes se classent comme prévu, et
l'arithmétique du polling et des frais.
\bfield{Pour aller plus loin} Des écarts de prix qui se referment~; plusieurs frontières à capacité partagée~; les fenêtres d'enchères.
\end{build}

\begin{omsources}
\item ENTSO-E, Single Intraday Coupling; ECCo SP (Energy Communication Platform).
\item Nord Pool, Intraday Market Regulations (juin 2024).
\item Betfair Exchange API documentation (Market Data Request Limits, Best Practice Guide) et articles d'assistance du Developer Program.
\end{omsources}

\section{Exercices}

\begin{exercise}[$\star$]\label{exo:nw:power-commodities-and-betting-connectivity:1}
Nommer les trois modules du système commun du couplage infrajournalier unique.
\end{exercise}

\begin{exercise}[$\star$]\label{exo:nw:power-commodities-and-betting-connectivity:2}
Expliquer comment deux ordres infrajournaliers au même prix sont classés sur Nord Pool.
\end{exercise}

\begin{exercise}[$\star$]\label{exo:nw:power-commodities-and-betting-connectivity:3}
Dire combien de marchés tiennent dans une requête de meilleurs cours vendeurs sur la bourse de paris, et pourquoi.
\end{exercise}

\begin{exercise}[$\star\star$]\label{exo:nw:power-commodities-and-betting-connectivity:4}
En s'appuyant sur la \cref{prop:nw:power-commodities-and-betting-connectivity:scarce}, expliquer pourquoi la firme conserve 95.4~\% des courses à
\qty{10}{\milli\second} quand la capacité restante vaut trois ordres, mais seulement 4.5~\% quand elle vaut un ordre~?
\end{exercise}

\begin{exercise}[$\star\star$]\label{exo:nw:power-commodities-and-betting-connectivity:5}
Expliquer pourquoi le marché continu cesse d'allouer la capacité inter-zones pendant une enchère infrajournalière sur la même frontière.
\end{exercise}

\begin{exercise}[$\star\star$]\label{exo:nw:power-commodities-and-betting-connectivity:6}
Un bot place et annule un pari 3\,000 fois par heure, et 500 de ses annulations échouent. Calculer le nombre de transactions qu'il compte.
\end{exercise}

\begin{exercise}[$\star\star\star$]\label{exo:nw:power-commodities-and-betting-connectivity:7}
\emph{Programmation.} Avec \texttt{nw\_power.share}, calculer la latence médiane qui donne à la firme la moitié des courses à un seul créneau.
\end{exercise}

\begin{exercise}[$\star\star\star$]\label{exo:nw:power-commodities-and-betting-connectivity:8}
\emph{Trouver la faille.} «~Le marché intrajournalier de l'électricité se négocie sur des durées de quelques minutes à quelques heures, donc la latence n'a pas d'importance.~»
\end{exercise}

\section{Problème~: Le dernier mégawatt à travers la frontière}

\begin{problem}[{Problème du week-end --- latence et capacité transfrontalière rare}]\label{pb:nw:power-commodities-and-betting-connectivity:1}
Un desk d'électricité négocie en intrajournalier à travers une frontière dont la capacité est souvent rare. Il affronte cinq rivaux pour la dernière capacité, comme dans le modèle du
chapitre, et valorise le spread à 12~EUR par mégawattheure en moyenne.

\medskip\noindent\textbf{Partie I --- Les règles.}
\begin{enumerate}
\item Comment la capacité transfrontalière est-elle allouée dans le trading continu intrajournalier~?
\item Que se passe-t-il pour la capacité lorsque deux ordres se croisent à travers la frontière~?
\item Quand le trading inter-zones s'arrête-t-il pour une période de livraison, et quand cela pourrait-il changer~?
\item Que se passe-t-il pendant une enchère intrajournalière sur la frontière~?
\end{enumerate}

\medskip\noindent\textbf{Partie II --- La course.}
\begin{enumerate}[resume]
\item Quelle part des courses à un seul créneau le desk remporte-t-il à 2 et à \qty{10}{\milli\second}~?
\item À quelle latence en remporte-t-il la moitié~?
\item Que change-t-on avec la capacité de trois ordres restante~?
\item Pourquoi le spread du rival le plus rapide compte-t-il~?
\end{enumerate}

\medskip\noindent\textbf{Partie III --- La valeur.}
\begin{enumerate}[resume]
\item Que espère le desk par mégawattheure proposé à 2 et à \qty{10}{\milli\second}~?
\item Pourquoi le deuxième participant le plus rapide n'obtient-il presque rien~?
\item Sur quelles frontières et à quelles heures la latence est-elle rentable~?
\item Que ignore le desk, et comment le découvrirait-il~?
\end{enumerate}

\medskip\noindent\textbf{Partie IV --- Le verdict.}
\begin{enumerate}[resume]
\item Énoncer le \emph{résultat nommé}~: la probabilité d'obtenir une capacité transfrontalière rare en fonction de la latence jusqu'au
  \omterm{def:nw:power-commodities-and-betting-connectivity:sob}{carnet d'ordres partagé}, et sa valeur par mégawattheure.
\item Quelles entrées sont publiées et lesquelles sont supposées~?
\item Par quel système de négociation local le desk doit-il passer~?
\item Que changerait un délai commun dans le module de capacité~?
\item Que doit mesurer le desk~?
\item Comment les programmes s'inscrivent-ils dans le même plan~?
\item Où cela se place-t-il dans le plan du chapitre~29~?
\item En une phrase~: quand un desk d'électricité intrajournalier doit-il être rapide~?
\end{enumerate}
\end{problem}

\section{Questions d'entretien}

\begin{interviewq}[$\star$ \iqroles{développeur}]\label{iq:nw:power-commodities-and-betting-connectivity:1}
Qu'est-ce que le couplage intrajournalier unique, et comment un ordre d'un pays rencontre-t-il un ordre d'un autre~?
\end{interviewq}

\begin{interviewq}[$\star\star$ \iqroles{trader}]\label{iq:nw:power-commodities-and-betting-connectivity:2}
Pourquoi quelques millisecondes peuvent-elles décider qui obtient la capacité transfrontalière en électricité intrajournalière~?
\end{interviewq}

\begin{interviewq}[$\star\star$ \iqroles{développeur}]\label{iq:nw:power-commodities-and-betting-connectivity:3}
Comment vous assureriez-vous que chaque transaction intrajournalière finisse dans un programme nommé avant l'échéance~?
\end{interviewq}

\begin{interviewq}[$\star\star$ \iqroles{développeur}]\label{iq:nw:power-commodities-and-betting-connectivity:4}
Faut-il interroger ou diffuser les cours d'une bourse de paris~? Lequel, et pourquoi~?
\end{interviewq}

\begin{interviewq}[$\star\star$ \iqroles{chercheur}]\label{iq:nw:power-commodities-and-betting-connectivity:5}
Comment estimeriez-vous la valeur d'être plus rapide vers le \omterm{def:nw:power-commodities-and-betting-connectivity:sob}{carnet d'ordres partagé}~?
\end{interviewq}

\begin{interviewq}[$\star\star\star$ \iqroles{développeur, chercheur}]\label{iq:nw:power-commodities-and-betting-connectivity:6}
Concevez la connectivité d'un desk de trading d'électricité européen qui négocie des contrats à terme, de l'intrajournalier et des programmes, et qui gère ses positions.
\end{interviewq}
